\subsection{Definitie en eerste eigenschappen}

\begin{lemma}[8.1.2]
    Zij $(V, \langle \cdot, \cdot \rangle)$ een inproduct-ruimte over $K = \mathbb{R}$ of $\mathbb{C}$. Dan gelden volgende eigenschappen, voor alle $u, v, w \in V$ en alle $\lambda, \mu \in K$.
    \begin{enumerate}
        \item[(i)] [\emph{bi-additiviteit}] $\langle u + v, w \rangle = \langle u, w \rangle + \langle v, w \rangle$ en $\langle u, v + w \rangle = \langle u, v \rangle + \langle u, w \rangle$;
        \item[(ii)] [\emph{sesqui-lineariteit}] $\langle \lambda v, \mu w \rangle = \lambda \overline{\mu} \cdot \langle v, w \rangle$;
        \item[(iii)] [\emph{homogeniteit van de norm}] $\|\lambda v\| = |\lambda| \cdot \|v\|$.
    \end{enumerate}
    Bovendien is $\langle \cdot, \cdot \rangle$ niet-ontaard, i.e. als $u \in V$ zodanig is dat $\langle u, v \rangle = 0$ voor alle $v \in V$, dan is noodzakelijk $u = 0$.
\end{lemma}
\begin{proof}(i)
    \(\langle u + v, w \rangle = \langle u, w \rangle + \langle v, w \rangle\) geldt uit (definitie 8.1.1) lineair in het eerste argument met \(a=b=1\).

    En 
    \(
        \langle u, v + w \rangle 
        = \overline{\langle v + w, u \rangle} 
        = \overline{\langle v , u \rangle} + \overline{\langle w, u \rangle}
        =\langle u, v \rangle + \langle u, w \rangle
    \)
\end{proof}
\begin{proof}(ii)
    \(
        \langle \lambda v, \mu w \rangle 
        = \lambda \langle  v, \mu w \rangle
        = \lambda \overline{\langle \mu w ,v\rangle}
        = \lambda \overline{\mu} \overline{\langle  w ,v\rangle}
        = \lambda \overline{\mu} \cdot \langle v, w \rangle
    \)
\end{proof}
\begin{proof}(iii)
    \(\|\lambda v\| 
    = \sqrt{\langle \lambda v, \lambda v \rangle} 
    = \sqrt{\lambda \overline{\lambda} \cdot \langle v, v \rangle}
    = \sqrt{\lambda \overline{\lambda}} \sqrt{ \langle v, v \rangle}
    = |\lambda| \cdot \|v\|\)
\end{proof}
Om aan te tonen dat  $\langle \cdot, \cdot \rangle$ niet-ontaard is, kiezen we \(u = v\).

\(\langle v, v \rangle = 0\).

Omdat \(\langle v, v \rangle \ngtr 0\) moet volgens (definitie 8.1.1) het positief definiet zijn, \(v=0\).


\begin{lemma}[8.2.2]
    Het orthogonaal complement $W^\perp$ van een deelruimte $W \le V$ is een deelruimte van $V$. Er geldt dat $W \cap W^\perp = \{0\}$.
\end{lemma}
\begin{proof}
    Per definitie van het orthogonaal complement is \(V = W \oplus W^\perp\).

    We bewijzen dat \(W^\perp\) een deelruimte is:
    Stel \(v_1,v_2 \in W^\perp\) en \(\lambda_1, \lambda_2 \in K\).

    \(\langle \lambda_1 v_1 + \lambda_2 v_2,w\rangle = \lambda_1 \langle  v_1  ,w\rangle + \lambda_2  \langle v_2,w\rangle = 0\).

    Stel \(v \in W \cap W^\perp\), dan \(\langle v,v \rangle = 0 \implies v = 0. \)
\end{proof}

\begin{lemma}[8.2.5]
    Zij $K = \mathbb{R}$ of $\mathbb{C}$, en zij $(V, \langle \cdot, \cdot \rangle)$ een eindig-dimensionale inproduct-ruimte over $K$. Voor elke deelruimte $W \le V$ is het orthogonaal complement $W^\perp$ in $V$ een complementaire ruimte,
    i.e. $W \oplus W^\perp = V$.
    In het bijzonder is $\dim(W^\perp) = \dim V - \dim W$.
\end{lemma}
\begin{proof}
    Uit (lemma 8.2.2) volgt  \(W \cap W^\perp = \{0\}\).
    
    Nu moeten we nog \(W + W^\perp = V\) aantonen.
    
    Als \(\{b_1, \dots, b_d\}\) een orthogonale basis voor \(W\) en \(v\in V\):
     \[v=\sum_{i=1}^{k} \langle v,b_i\rangle b_i + \left(v - \sum_{i=1}^{k} \langle v,b_i\rangle b_i \right).\]
    De eerste term is een element van \(W\). 
    Neem \(b_j \in \{b_1,\dots,b_d\}\) willekeurig, dan is
    \[
    \langle v - \sum_{i}^{d} \langle v,b_j \rangle b_i, b_j \rangle 
    = \langle v,b_j\rangle - \sum_{i}^{d} \langle v,b_i \rangle \langle b_i,b_j \rangle
    = \langle v,b_j \rangle - \langle v,b_j \rangle = 0.
    \]
\end{proof}

\begin{gevolg}[8.2.6]
    Zij $K = \mathbb{R}$ of $\mathbb{C}$, en zij $(V, \langle \cdot, \cdot \rangle)$ een eindig-dimensionale inproduct-ruimte over $K$. Voor elke deelruimte $W \le V$ is $(W^\perp)^\perp = W$.
\end{gevolg}
\begin{proof}
    \(\dim (W^\perp) = \dim(V)-\dim(W) \)

    \(\dim((W^\perp)^\perp) = \dim(V) -(\dim(V)-\dim(W)) = \dim(W)\)
\end{proof}

\begin{lemma}[8.2.9]
    Zij $K = \mathbb{R}$ of $\mathbb{C}$, en zij $(V, \langle \cdot, \cdot \rangle)$ een eindig-dimensionale inproduct-ruimte over $K$. Beschouw een deelruimte $W \le V$.
    \begin{enumerate}
        \item[(i)] Voor alle $w \in W$ is $\operatorname{proj}_W(w) = w$; voor alle $u \in W^\perp$ is $\operatorname{proj}_W(u) = 0$. In het bijzonder is de orthogonale projectie een projectie-operator in de betekenis van Definitie 3.1.7.
        \item[(ii)] Zij $\{w_1, \dots, w_k\}$ een orthonormale basis van $W$. Dan is
        \[ \operatorname{proj}_W(v) = \sum_{j=1}^k \langle v, w_j \rangle w_j \quad \forall v \in V.\]
    \end{enumerate}
\end{lemma}

\begin{proof}(i)
    Volgt uit \(w =w + 0 \in W \oplus W^\perp \) en \(u = 0 + u \in W \oplus W^\perp\).
\end{proof}
\begin{proof}(ii)
    Net als in het bewijs van gevolg 8.2.5 is
    \[v=\sum_{i=1}^{k} \langle v,w_i\rangle w_i + \left(v - \sum_{i=1}^{k} \langle v,w_i\rangle w_i \right) \in W \oplus W^\perp.\]
\end{proof}



